\documentclass[11pt]{letter}
\usepackage[margin=1in]{geometry}
\usepackage{amsmath,amssymb}

\signature{Raeez Lorgat\\
Perimeter Institute for Theoretical Physics\\
31 Caroline Street North\\
Waterloo, ON N2L 2Y5, Canada\\
\texttt{rlorgat@perimeterinstitute.ca}}

\date{\today}

\begin{document}

\begin{letter}{The Editors\\
\emph{Letters in Mathematical Physics}}

\opening{Dear Editors,}

This standalone note is not submission-ready. It records only the
primary-state scalar Ward factor for the Virasoro collision-residue
connection, not the full completed Virasoro spectral $R$-operator.
On a highest-weight line of conformal weight~$h$ the scalar factor is
\[
R^{\mathrm{prim}}_{h,c}(z)
= z^{2h}\exp\!\left(-\frac{c}{4z^2}\right).
\]
The derivation proceeds from the collision residue
$r^{\mathrm{coll}}(z) = (c/2)/z^3 + 2T/z$ of the Virasoro OPE via
$d\log$ absorption (reducing pole orders by one) and path-ordered
exponentiation on the commutative primary sector.

Only even powers of $1/z$ appear after factoring out $z^{2h}$. The
scalar coefficient $R_2=-c/4$ is a normalization check for the
highest-weight quotient; it is not by itself a proof of class~$M$
membership or of $A_\infty$ non-formality. Those statements require
the completed Virasoro obstruction tower and the wheel-survival
hypotheses in the monograph.

\closing{Sincerely,}

\end{letter}
\end{document}
